\documentclass[10pt,a4paper]{article}

% Kodowanie, język i typografia
\usepackage[T1]{fontenc}
\usepackage[utf8]{inputenc}
\usepackage[polish]{babel}
\usepackage{lmodern}
\usepackage{microtype}

% Układ strony i odstępy
\usepackage[a4paper,margin=2.5cm]{geometry}
\usepackage{setspace}
\setstretch{1.25}
\usepackage{parskip}
\usepackage{enumitem}

% Tabele, matematyka i odnośniki
\usepackage{amsmath}
\usepackage{mathtools}
\usepackage{booktabs}
\usepackage{tabularx}
\usepackage{array}
\usepackage{hyperref}
\usepackage{xurl}

\hypersetup{
	pdftitle={Wpływ cech przestrzennych i strategii walidacji na dokładność predykcji cen najmu mieszkań z wykorzystaniem modelu CatBoost},
	pdfauthor={Ivan Kokhman, Błażej Krasucki},
	hidelinks
}

\urlstyle{same}

\setlist[enumerate]{
	leftmargin=*,
	itemsep=0.25em,
	topsep=0.25em
}

\setlist[itemize]{
	leftmargin=*,
	itemsep=0.25em,
	topsep=0.25em
}

\renewcommand{\arraystretch}{1.2}

% Tytuł dokumentu
\title{
	\textbf{
		Wpływ cech przestrzennych i strategii walidacji
		na dokładność predykcji cen najmu mieszkań
		z wykorzystaniem modelu CatBoost
	}
	\\[0.6em]
}

\author{
	Ivan Kokhman (104432)\\
	Błażej Krasucki (103885)\\[0.8em]
	{\normalsize Politechnika Opolska -- kierunek Informatyka}\\[0.3em]
	{\normalsize Przedmiot: Rozpoznawanie wzorców}\\[0.8em]
	{\normalsize October 2026}
}

\date{}

\begin{document}
	
	\maketitle
	
	
	%------------------------------------------------
	% STRESZCZENIE
	%------------------------------------------------
	
	\begin{abstract}
		Przedmiotem planowanej pracy jest ocena wpływu cech przestrzennych
		oraz wyboru strategii walidacji na jakość predykcji cen najmu mieszkań.
		Badanie będzie realizowane na danych ogłoszeniowych dotyczących mieszkań
		w polskich miastach, zebranych w ramach projektu SmartRent AI.
		Porównane zostaną regresja liniowa jako model odniesienia i model CatBoost,
		trenowane zarówno na podstawowych cechach lokalu, jak i na zbiorze
		rozszerzonym o informacje przestrzenne.
		Zastosowane zostaną losowy podział danych oraz walidacja grupowa
		uwzględniająca lokalizację.
		Wyniki zostaną ocenione za pomocą MAE, RMSE i współczynnika $R^2$.
	\end{abstract}
	
	\noindent

	\noindent
	\textbf{Repozytorium projektu:}
	\url{https://github.com/voloshka-eth/smart-rent}
	
	%------------------------------------------------
	% 1. TEMAT, CEL I UZASADNIENIE
	%------------------------------------------------
	\newpage
	\section{Temat, cel i uzasadnienie}
	
	Na cenę najmu mieszkania wpływają zarówno parametry lokalu,
	jak i jego położenie. Losowy podział danych może umieszczać
	podobne, sąsiednie oferty jednocześnie w zbiorach uczącym
	i testowym, prowadząc do optymistycznej oceny predykcji
	\cite{roberts2017}.
	
	\noindent
	\textbf{Cel pracy:}
	ocenić wpływ cech przestrzennych i strategii walidacji
	na błąd predykcji ceny najmu.
	
	\subsection{Pytania badawcze i hipotezy}
	
	W ramach badania zostaną zweryfikowane następujące hipotezy:
	
	\begin{itemize}
		\item CatBoost osiągnie niższy MAE niż regresja liniowa.
		\item Dodanie cech przestrzennych zmniejszy błąd predykcji.
		\item Walidacja przestrzenna wykaże większy błąd niż podział losowy.
	\end{itemize}
	
	%------------------------------------------------
	% 2. OPIS ZADANIA I DANE WEJŚCIOWE
	%------------------------------------------------
	
	\section{Opis zadania i dane wejściowe}
	
	Wykorzystane zostaną dane z projektu SmartRent AI:
	wstępnie około 18 tys. ogłoszeń z dziewięciu polskich miast.
	Po oczyszczeniu zbioru model otrzyma cechy mieszkania
	(miasto, dzielnica, powierzchnia, liczba pokoi, piętro
	i udogodnienia) oraz -- w wariancie rozszerzonym --
	odległości do centrum, transportu i wybranych usług.
	
	Przygotowanie danych obejmie deduplikację, odfiltrowanie
	błędnych ofert i ocenę jakości geolokalizacji.
	Celem predykcji jest cena ofertowa najmu w PLN miesięcznie,
	bez opłat administracyjnych.
	Zmienna \texttt{price\_per\_m2} zostanie wykluczona
	z predyktorów z uwagi na ryzyko przecieku informacji.
	
	%------------------------------------------------
	% 3. PLANOWANA METODYKA
	%------------------------------------------------
	
	\section{Planowana metodyka}
	
	Porównane zostaną regresja liniowa (model bazowy) oraz
	\texttt{CatBoostRegressor} \cite{prokhorenkova2018}.
	Każdy model zostanie uruchomiony z podstawowym i przestrzennie
	rozszerzonym zestawem cech, co da łącznie cztery konfiguracje
	eksperymentu.
	
	W ramach porównania zostaną zastosowane dwie strategie walidacji:
	
	\begin{enumerate}
		\item losowa walidacja krzyżowa;
		\item walidacja grupowa, rozdzielająca obszary
		miasto--dzielnica pomiędzy foldy.
	\end{enumerate}
	
	W obu strategiach wykorzystane będą identyczne konfiguracje modeli.
	Przygotowanie cech nastąpi wyłącznie na danych treningowych,
	aby ograniczyć ryzyko przecieku informacji.
	
	Wyniki zostaną przedstawione za pomocą metryk MAE, RMSE i $R^2$,
	również w rozbiciu na poszczególne miasta.
	
	%------------------------------------------------
	% 4. PROPONOWANE CHECKPOINTY
	%------------------------------------------------
	
	\newpage
	\section{Proponowane checkpointy}
	
	\begin{table}[htbp]
		\centering
		\small
		\caption{Harmonogram realizacji projektu}
		\label{tab:checkpointy}
		
		\begin{tabularx}{\textwidth}{
				@{}
				p{2.1cm}
				p{4.2cm}
				X
				@{}
			}
			\toprule
			\textbf{Termin}
			&
			\textbf{Etap}
			&
			\textbf{Oczekiwany rezultat / dowód wykonania}
			\\
			\midrule
			
			27.10.2026
			&
			Checkpoint 1: dane i plan eksperymentu
			&
			Zweryfikowany zbiór danych, raport jakości geolokalizacji,
			lista cech i ustalony schemat podziałów.
			\\
			
			24.11.2026
			&
			Checkpoint 2: modele i walidacja
			&
			Uruchomione modele bazowe i CatBoost, porównanie wariantów
			cech oraz tabela MAE/RMSE/$R^2$ dla obu schematów walidacji.
			\\
			
			12.01.2027
			&
			Checkpoint 3: analiza wyników i artykuł
			&
			Interpretacja hipotez, wykresy i tabele, analiza ograniczeń
			oraz kompletna wersja miniartykułu.
			\\
			
			\bottomrule
		\end{tabularx}
	\end{table}
	
	%------------------------------------------------
	% 5. PLANOWANE REZULTATY
	%------------------------------------------------
	
	\section{Planowane rezultaty i kryteria zakończenia}
	
	Rezultatem projektu będzie artykuł zawierający tabelę metryk,
	opis eksperymentów, porównanie strategii walidacji oraz weryfikację
	postawionych hipotez.
	
	Repozytorium będzie zawierać kod i parametry pozwalające odtworzyć
	wyniki. W opracowaniu zostaną opisane zarówno wyniki potwierdzające
	hipotezy, jak i wyniki im przeczące.
	
	%------------------------------------------------
	% BIBLIOGRAFIA
	%------------------------------------------------
	
	\begin{thebibliography}{99}
		
				\bibitem{roberts2017}
		D. R. Roberts et al.,
		\newblock ``Cross-validation strategies for data with temporal,
		spatial, hierarchical, or phylogenetic structure,''
		\newblock \emph{Ecography}, vol.~40, no.~8, pp.~913--929, 2017.
		\newblock DOI:
		\href{https://doi.org/10.1111/ecog.02881}
		{10.1111/ecog.02881}.
		
		\bibitem{prokhorenkova2018}
		L. Prokhorenkova et al.,
		\newblock ``CatBoost: Unbiased Boosting with Categorical Features,''
		\newblock \emph{Advances in Neural Information Processing Systems},
		vol.~31, 2018.
		

		
	\end{thebibliography}
	
\end{document}